% Authoritative LaTeX chapter. Edit this file for future revisions.
% Initially migrated 2026-09-11; no Markdown import occurs during PDF builds.
\chapter{출력 포맷과 결과 해석}
\label{chap:05_output_formats}
매뉴얼 목차 · \hyperref[chap:01_concepts_geometry]{기본 개념과 RAA}

이 장은 Git 커밋 \texttt{54861c68e35ec5aa\allowbreak{}a3938fad53e93dc2\allowbreak{}dfa7eb1d}의 C 실행 프로그램이 실제로 쓰는 헤더와 출력문을 기준으로 한다. 과거 파일에 같은 뜻의 다른 열 이름이 있을 수 있으므로 열의 위치보다 \textbf{실제 헤더 이름}으로 읽는다.

\section{실행 모드별 결과물}
\label{sec:auto:05_output_formats:1}

{\TableFont
\begin{longtable}{@{}>{\raggedright\arraybackslash}p{\dimexpr 0.3000\linewidth-5.33pt\relax}>{\raggedright\arraybackslash}p{\dimexpr 0.2300\linewidth-5.33pt\relax}>{\raggedright\arraybackslash}p{\dimexpr 0.4700\linewidth-5.33pt\relax}@{}}
\toprule
\textbf{모드} & \textbf{주요 출력} & \textbf{쓰는 위치} \\
\midrule
\endfirsthead
\toprule
\textbf{모드} & \textbf{주요 출력} & \textbf{쓰는 위치} \\
\midrule
\endhead
\bottomrule
\endfoot
단일 기하: \texttt{-\allowbreak{}-\allowbreak{}vza}와 \texttt{-\allowbreak{}-\allowbreak{}raa} 모두 지정 & 처음에 I/Q/U 숫자 3개, 이어서 \texttt{key=\allowbreak{}value}가 나오는 한 줄 & 표준 출력(stdout); 필요하면 파일로 리다이렉션 \\
비해양 각도 격자: \texttt{black}, \texttt{flat}, \texttt{black\_\allowbreak{}fresnel\_\allowbreak{}ocean} & 대기·표면 LUT CSV & \texttt{-\allowbreak{}-\allowbreak{}output-\allowbreak{}full-\allowbreak{}grid FILE.\allowbreak{}csv}; 생략 시 현재 디렉터리의 \texttt{result<N>.\allowbreak{}csv} \\
해양 각도 격자: \texttt{ocean} & \texttt{Rrs}, \texttt{rrs}, TOA, 조도·투과·IOP를 포함한 CSV & 같은 파일 옵션 \\
\texttt{-\allowbreak{}-\allowbreak{}batch CSV -\allowbreak{}-\allowbreak{}output FILE.\allowbreak{}csv} & 입력 순서의 지정 기하 결과 CSV & \texttt{-\allowbreak{}-\allowbreak{}output} \\
\texttt{-\allowbreak{}-\allowbreak{}batch}와 격자 출력을 함께 요청 & 요약 CSV와 행별 \texttt{case\_\allowbreak{}<case\_\allowbreak{}id>.\allowbreak{}csv} & \texttt{-\allowbreak{}-\allowbreak{}output}과 \texttt{-\allowbreak{}-\allowbreak{}output-\allowbreak{}full-\allowbreak{}grid}로 지정한 디렉터리 \\
\texttt{-\allowbreak{}-\allowbreak{}batch-\allowbreak{}full-\allowbreak{}grid JOBS.\allowbreak{}csv} & job마다 해양 격자 CSV & 각 job의 \texttt{out} 열 \\
\texttt{-\allowbreak{}-\allowbreak{}ccrr-\allowbreak{}compare} & OCRT 값·참조값·상대오차 CSV & \texttt{-\allowbreak{}-\allowbreak{}ccrr-\allowbreak{}output} 또는 해당 비교의 \texttt{-\allowbreak{}-\allowbreak{}output} \\
\end{longtable}
}

공개 C 실행 프로그램의 결과 파일 형식은 CSV 또는 단일 기하 텍스트다. JSON·NetCDF·HDF5는 이 출력기의 기본 형식이 아니다. 저장소의 검증 JSON이나 도구가 만든 통합 파일은 각 도구의 별도 산출물이다.

CSV는 첫 줄에 헤더, 이후 한 줄에 한 결과를 쓴다. 구분자는 쉼표, 소수점은 \texttt{.\allowbreak{}}, 작은 값은 \texttt{1.\allowbreak{}234e-\allowbreak{}05}와 같은 과학적 표기다. \texttt{nan}은 계산 불가·정의되지 않음일 수 있다. 결측값을 0으로 바꾸기 전에 해당 열과 유효성 플래그를 확인한다. 파일 출력과 별개로 진단·진행 메시지는 대부분 stderr에 나온다. 배치 요약은 stdout에도 나오므로 배치 stdout 전체를 CSV로 취급하지 않는다.

\textbf{열 이름의 대소문자를 보존한다.} 해양 CSV의 \texttt{Rrs\_\allowbreak{}I}와 \texttt{rrs\_\allowbreak{}I}는 서로 다른 물리량이다. 열 이름을 일괄 소문자로 바꾸지 않는다. 읽기 도구 선택과 Python 예시는 부록~\ref{sec:example_csv_reader}를 참고한다.

격자 CSV는 \textbf{VZA가 바깥쪽, RAA가 안쪽 루프}다. 같은 VZA에서 RAA가 증가한 뒤 다음 VZA로 넘어간다. RAA \texttt{360\textdegree{}} 중복 끝점은 없다. 간격이 360이나 VZA 상한을 정확히 나누지 않으면 마지막 출력 좌표가 요청 상한과 다를 수 있으므로 파일에 기록된 좌표를 사용한다. 격자의 각도 간격과 SOS 적분용 Gauss 노드 수는 서로 다른 설정이다.

\Needspace{0.60\textheight}
\section{단위와 입사광 정규화}
\label{sec:auto:05_output_formats:2}
그림~\ref{fig:output-levels}에서 먼저 출력 높이와 각 비율의 분모를 구분한다.
% Height conventions and normalization, schematic rather than a ray solver.
\begin{figure}[H]
\centering
\begin{tikzpicture}[x=1cm,y=1cm,>=Latex,
 every node/.style={font=\sffamily\fontsize{9}{11}\selectfont},
 quantity/.style={draw=ocrtLine,rounded corners=2pt,fill=white,text width=6.5cm,align=left,inner sep=6pt}]
\fill[ocrtTeal!9] (0,0) rectangle (6.7,1.7);
\fill[blue!3] (0,1.7) rectangle (6.7,5.1);
\draw[ocrtNavy,thick] (0,1.7)--(6.7,1.7);
\draw[ocrtGray,dashed] (0,5.1)--(6.7,5.1);
\node[anchor=west] at (0.1,5.35) {TOA};
\node[anchor=west,text=ocrtGray] at (0.1,4.5) {\FigLang{대기}{Atmosphere}};
\node[anchor=west,text=ocrtTeal] at (0.1,0.25) {\FigLang{해수}{Seawater}};
\node[anchor=east] at (6.6,1.97) {$0^+$};
\node[anchor=east] at (6.6,1.40) {$0^-$};
\draw[->,very thick,orange!85!black] (1.1,3.6)--(1.1,1.9);
\node[anchor=west,align=left] at (1.35,3.05) {$E_d(0^+)$\\\FigLang{하향 조도}{Downward irradiance}};
\draw[->,very thick,orange!85!black] (1.1,1.5)--(1.1,0.6);
\node[anchor=west] at (1.35,0.98) {$E_d(0^-)$};
\draw[->,very thick,ocrtTeal] (4.25,0.3)--(4.65,1.5);
\node[anchor=west] at (4.65,0.55) {$L_u(0^-)$};
\draw[->,very thick,ocrtTeal] (4.7,1.9)--(6.2,4.65);
\node[anchor=east,align=right] at (5.2,3.85) {$L_u(0^+)$\\\FigLang{상향 휘도}{Upward radiance}};
\draw[->,very thick,ocrtNavy] (6.2,4.8)--(6.65,5.62);
\node[anchor=east] at (6.1,5.75) {$L_{\rm TOA}$};
\node[quantity,anchor=north west] at (7.5,5.9)
 {$\displaystyle \rho_X=\frac{\pi L_{{\rm TOA},X}}{F_0\cos({\rm SZA})}$\\[4pt]
  \FigLang{TOA 반사도: 무차원}{TOA reflectance: dimensionless}};
\node[quantity,anchor=north west] at (7.5,3.83)
 {$\displaystyle Rrs_X=\frac{L_{u,X}(0^+)}{E_d(0^+)}$\\[4pt]
  \FigLang{수면 위: sr$^{-1}$}{Above water: sr$^{-1}$}};
\node[quantity,anchor=north west] at (7.5,1.76)
 {$\displaystyle rrs_X=\frac{L_{u,X}(0^-)}{E_d(0^-)}$\\[4pt]
  \FigLang{수면 아래: sr$^{-1}$}{Below water: sr$^{-1}$}};
\node[anchor=north west,text width=14.6cm,align=left,font=\sffamily\footnotesize] at (0,-0.35)
 {$X\in\{I,Q,U\}$\qquad
  \FigLang{$F_0$: 태양광에 수직인 면의 입사량; 수평면에서는 $F_0\cos({\rm SZA})$}{$F_0$: incident irradiance normal to the solar beam; a horizontal plane receives $F_0\cos({\rm SZA})$}};
\end{tikzpicture}
\caption{\FigLang{출력 위치와 정규화의 구분. $0^+$와 $0^-$는 수면을 사이에 둔 서로 다른 위치다. 대문자 \texttt{Rrs}와 소문자 \texttt{rrs}는 교환할 수 없고, 둘에 $\pi$를 곱해도 대기 경로를 포함한 TOA \texttt{rho}가 되지 않는다. 화살표는 물리량의 방향과 평가 위치를 설명하는 도식이며 모든 산란 경로를 그린 것은 아니다.}{Output levels and their normalizations. $0^+$ and $0^-$ are distinct positions on opposite sides of the interface. Uppercase \texttt{Rrs} and lowercase \texttt{rrs} are not interchangeable. Multiplying either by $\pi$ does not produce TOA \texttt{rho}, which includes the atmosphere. Arrows identify directions and sampling levels schematically; they do not enumerate every scattering path.}}
\label{fig:output-levels}
\end{figure}

\FloatBarrier
\texttt{F0}를 태양광 진행 방향에 수직인 면의 TOA 분광 입사량, \texttt{mu\_\allowbreak{}s=\allowbreak{}cos(SZA)}라 하면 수평면의 입사량은 \texttt{F0*mu\_\allowbreak{}s}다.


\begin{ManualCode}
rho_X = pi * L_TOA_X / (F0 * mu_s)          X = I, Q, U
Rrs_X = Lu_X(0+) / Ed(0+)                  [sr^-1]
rrs_X = Lu_X(0-) / Ed(0-)                  [sr^-1]
BRF0plus_X  = pi * Rrs_X                   [무차원]
BRF0minus_X = pi * rrs_X                   [무차원]
\end{ManualCode}

대기 SOS는 \texttt{F0=\allowbreak{}\ensuremath{\pi}}를 소스 정규화에 포함하므로 내부 TOA intensity에서 \texttt{rho=\allowbreak{}I\_\allowbreak{}TOA/\allowbreak{}mu\_\allowbreak{}s}로 변환한다. 출력 \texttt{rho}에 \ensuremath{\pi}나 \texttt{1/\allowbreak{}mu\_\allowbreak{}s}를 다시 곱하지 않는다.

수중 \texttt{Lu/\allowbreak{}Qu/\allowbreak{}Uu}, \texttt{Ed/\allowbreak{}Eu}는 각각 복사휘도·복사조도 차원의 값이지만 \textbf{실제 날짜·센서의 태양 스펙트럼을 읽어서 절대 복사량으로 만든 값은 아니다.} 현재 C 소스의 입사광 스케일은 다음과 같다.


{\TableFont
\begin{longtable}{@{}>{\raggedright\arraybackslash}p{\dimexpr 0.3400\linewidth-4.00pt\relax}>{\raggedright\arraybackslash}p{\dimexpr 0.6600\linewidth-4.00pt\relax}@{}}
\toprule
\textbf{물 계산 경로} & \textbf{내부 \texttt{F\_\allowbreak{}sun}} \\
\midrule
\endfirsthead
\toprule
\textbf{물 계산 경로} & \textbf{내부 \texttt{F\_\allowbreak{}sun}} \\
\midrule
\endhead
\bottomrule
\endfoot
원래 순수 해수 경로; OCRT/CCRR 구성성분을 모두 0으로 지정해 순수 해수로 전환한 경로 & \ensuremath{\pi} \\
0이 아닌 구성성분을 사용하는 OCRT/CCRR 경로 & 1 \\
직접 IOP 경로 & 1 \\
전용 CCRR 참조 비교 & 1 \\
\end{longtable}
}

따라서 서로 다른 경로의 원시 \texttt{Lu/\allowbreak{}Ed} 숫자를 그대로 비교하면 정규화 차이가 섞일 수 있다. \texttt{Rrs}, \texttt{rrs}, \texttt{rho}는 각자의 입사량으로 정규화한 양이다. 물의 원시 radiance·irradiance를 실제 \texttt{F0\_\allowbreak{}actual}에 맞추려면 사용한 \texttt{F\_\allowbreak{}sun}을 확인하고 동일한 물리 단위에서 \texttt{F0\_\allowbreak{}actual/\allowbreak{}F\_\allowbreak{}sun} 배율을 적용한다. TOA 절대 복사휘도는 정규화된 \texttt{rho}에서 \texttt{L\_\allowbreak{}TOA=\allowbreak{}rho*F0\_\allowbreak{}actual*cos(SZA)/\allowbreak{}\ensuremath{\pi}}로 복원한다. CLI의 과거 \texttt{-\allowbreak{}-\allowbreak{}f-\allowbreak{}sun} 옵션은 제거되었다.


{\TableFont
\begin{longtable}{@{}>{\raggedright\arraybackslash}p{\dimexpr 0.3400\linewidth-4.00pt\relax}>{\raggedright\arraybackslash}p{\dimexpr 0.6600\linewidth-4.00pt\relax}@{}}
\toprule
\textbf{수량} & \textbf{단위·해석} \\
\midrule
\endfirsthead
\toprule
\textbf{수량} & \textbf{단위·해석} \\
\midrule
\endhead
\bottomrule
\endfoot
\texttt{rho}, \texttt{BRF}, \texttt{omega}, AOD, \texttt{tau\_\allowbreak{}R}, \texttt{T\_\allowbreak{}*} & 무차원; 방향별 반사도·전달비는 항상 0--1에 갇힌 에너지 확률이 아님 \\
\texttt{Rrs}, \texttt{rrs} & sr\textsuperscript{-}\textsuperscript{1} \\
\texttt{Lu}, \texttt{Qu}, \texttt{Uu}, \texttt{TOA\_\allowbreak{}water\_\allowbreak{}signal\_\allowbreak{}I} & 위 입사광 스케일의 분광 복사휘도; 실제 스케일 부여 시 W m\textsuperscript{-}\textsuperscript{2} sr\textsuperscript{-}\textsuperscript{1} nm\textsuperscript{-}\textsuperscript{1} \\
\texttt{Ed}, \texttt{Eu} & 위 입사광 스케일의 분광 복사조도; 실제 스케일 부여 시 W m\textsuperscript{-}\textsuperscript{2} nm\textsuperscript{-}\textsuperscript{1} \\
\texttt{a}, \texttt{b}, \texttt{bb}, \texttt{Kd}, \texttt{Ku} & m\textsuperscript{-}\textsuperscript{1} \\
기하 & 도; \texttt{mu\_\allowbreak{}view}는 \texttt{cos(VZA)}라 무차원 \\
\texttt{wavelength\_\allowbreak{}nm}, \texttt{AOD\_\allowbreak{}ref\_\allowbreak{}nm} & nm \\
\texttt{wind\_\allowbreak{}speed\_\allowbreak{}ms} & m s\textsuperscript{-}\textsuperscript{1} \\
\end{longtable}
}

\section{선글린트와 TOA 성분을 구분한다}
\label{sec:auto:05_output_formats:3}
해양 \texttt{ocean} 단일 기하 출력은 직접 선글린트를 분리한 기본 TOA와 직접 선글린트를 더한 TOA를 함께 제공한다. \texttt{X}는 각각 \texttt{I/\allowbreak{}Q/\allowbreak{}U}를 뜻한다.


\begin{ManualCode}
TOA_rho_water_total_X = TOA_rho_water_direct_X + TOA_rho_water_sky_X
TOA_rho_X = TOA_rho_atm_X + TOA_rho_water_total_X
TOA_rho_glint_on_X = TOA_rho_X + TOA_rho_glint_direct_X
\end{ManualCode}


{\TableFont
\begin{longtable}{@{}>{\raggedright\arraybackslash}p{\dimexpr 0.3400\linewidth-4.00pt\relax}>{\raggedright\arraybackslash}p{\dimexpr 0.6600\linewidth-4.00pt\relax}@{}}
\toprule
\textbf{열·키} & \textbf{의미} \\
\midrule
\endfirsthead
\toprule
\textbf{열·키} & \textbf{의미} \\
\midrule
\endhead
\bottomrule
\endfoot
\texttt{TOA\_\allowbreak{}rho\_\allowbreak{}atm\_\allowbreak{}I/\allowbreak{}Q/\allowbreak{}U} & 결합 계산에서 별도 대기·표면 경로로 계산한 성분; 순수 Rayleigh만을 뜻하지 않음 \\
\texttt{TOA\_\allowbreak{}rho\_\allowbreak{}water\_\allowbreak{}direct\_\allowbreak{}I/\allowbreak{}Q/\allowbreak{}U} & 직접 태양광으로 시작한 수중 신호가 수면을 나와 대기를 거쳐 TOA에 도달한 성분 \\
\texttt{TOA\_\allowbreak{}rho\_\allowbreak{}water\_\allowbreak{}sky\_\allowbreak{}I/\allowbreak{}Q/\allowbreak{}U} & 대기 하향 확산광으로 시작한 수중 신호의 TOA 성분 \\
\texttt{TOA\_\allowbreak{}rho\_\allowbreak{}water\_\allowbreak{}total\_\allowbreak{}I/\allowbreak{}Q/\allowbreak{}U} & 위 두 물 기원 성분의 합 \\
\texttt{TOA\_\allowbreak{}rho\_\allowbreak{}glint\_\allowbreak{}direct\_\allowbreak{}I/\allowbreak{}Q/\allowbreak{}U} & 수면의 직접 태양광 정반사 항 \\
\texttt{TOA\_\allowbreak{}rho\_\allowbreak{}glint\_\allowbreak{}on\_\allowbreak{}I/\allowbreak{}Q/\allowbreak{}U} & 기본 TOA에 직접 선글린트를 더한 값 \\
\end{longtable}
}

\texttt{water\_\allowbreak{}direct}의 direct는 \textbf{입사광의 기원}이다. 그 물 신호가 물속이나 대기에서 한 번도 산란하지 않았다는 의미가 아니다. 아래 \texttt{Ed0plus\_\allowbreak{}direct}, \texttt{Ed0minus\_\allowbreak{}direct}의 direct는 \textbf{그 위치의 비산란 평행 태양광}이므로 구분한다.

해양 격자 CSV의 \texttt{TOA\_\allowbreak{}rho\_\allowbreak{}I/\allowbreak{}Q/\allowbreak{}U}도 이 해양 기본 TOA다. 그러나 현재 격자 출력기는 \texttt{TOA\_\allowbreak{}rho\_\allowbreak{}glint\_\allowbreak{}on\_\allowbreak{}*}와 위 TOA 성분 분해 열을 쓰지 않는다. 해당 성분이 필요한 모의자료는 단일 기하 출력의 이름 있는 키를 수집한다.

비해양 \texttt{black\_\allowbreak{}fresnel\_\allowbreak{}ocean}의 \texttt{rho\_\allowbreak{}I/\allowbreak{}Q/\allowbreak{}U}는 \texttt{-\allowbreak{}-\allowbreak{}decouple-\allowbreak{}sunglint} 설정에 따라 직접 선글린트 포함 여부가 달라진다. 이를 해양 CSV의 동일한 \texttt{rho} 개념과 혼동하지 않는다. 풍속 0의 완전 평탄면 정반사는 유한한 폭의 Cox--Munk 글린트와 다르며, 정확한 정반사 방향의 값은 바람이 있는 경우처럼 해석할 수 없다.

\textbf{현재 비해양 U 진단의 주의점:} \texttt{rt\_\allowbreak{}solver.\allowbreak{}c}는 비해양 총량에 직접 글린트 \texttt{-U\_\allowbreak{}kernel}을 더하지만, 비해양 \texttt{TOA\_\allowbreak{}rho\_\allowbreak{}glint\_\allowbreak{}direct\_\allowbreak{}U}에는 \texttt{+U\_\allowbreak{}kernel}을 기록한다. 해양 \texttt{ocean}의 같은 이름 키는 총량에 더하는 부호와 일치한다. 따라서 비해양 키를 그대로 총량에서 빼서 glint-free U를 만들지 않는다. 비해양 편광 분리는 \texttt{-\allowbreak{}-\allowbreak{}decouple-\allowbreak{}sunglint} 실행 결과로 확인한다. 이 문서는 해당 출력 부호 불일치를 기록하며 코드를 바꾸지는 않는다.

\section{단일 기하 텍스트}
\label{sec:auto:05_output_formats:4}
형태는 다음과 같다. 아래는 형식 예시이며 수치 예측값이 아니다.


\begin{ManualCode}
<rho_I> <rho_Q> <rho_U>  TOA_rho_I=<value> ... sza=<value> vza=<value> raa=<value> wl=<value> orders=<integer> conv=<0-or-1>
\end{ManualCode}

첫 세 숫자는 이름 있는 \texttt{TOA\_\allowbreak{}rho\_\allowbreak{}I/\allowbreak{}Q/\allowbreak{}U}와 같은 값이다. 새 분석 코드는 위치 인덱스보다 \texttt{key=\allowbreak{}value}를 파싱한다. \texttt{-\allowbreak{}-\allowbreak{}output-\allowbreak{}advanced}는 배치 CSV의 열을 늘리는 옵션이며 이 단일 기하 출력의 형식을 바꾸지 않는다.

\subsection{모든 표면에 공통인 키}
\label{sec:auto:05_output_formats:4.1}

{\TableFont
\begin{longtable}{@{}>{\raggedright\arraybackslash}p{\dimexpr 0.3400\linewidth-4.00pt\relax}>{\raggedright\arraybackslash}p{\dimexpr 0.6600\linewidth-4.00pt\relax}@{}}
\toprule
\textbf{키} & \textbf{의미} \\
\midrule
\endfirsthead
\toprule
\textbf{키} & \textbf{의미} \\
\midrule
\endhead
\bottomrule
\endfoot
\texttt{TOA\_\allowbreak{}rho\_\allowbreak{}I}, \texttt{TOA\_\allowbreak{}rho\_\allowbreak{}Q}, \texttt{TOA\_\allowbreak{}rho\_\allowbreak{}U} & TOA Stokes 방향 반사도; 글린트 포함 조건은 \ref{sec:auto:05_output_formats:3}절 \\
\texttt{TOA\_\allowbreak{}rho\_\allowbreak{}glint\_\allowbreak{}direct\_\allowbreak{}I/\allowbreak{}Q/\allowbreak{}U} & 직접 선글린트 진단; 비해양 U 부호 주의 \\
\texttt{T\_\allowbreak{}dir\_\allowbreak{}dn}, \texttt{T\_\allowbreak{}diff\_\allowbreak{}dn\_\allowbreak{}hemi}, \texttt{T\_\allowbreak{}total\_\allowbreak{}dn\_\allowbreak{}hemi}, \texttt{T\_\allowbreak{}diff\_\allowbreak{}dn\_\allowbreak{}dir} & 아래 하향 투과 정의 \\
\texttt{T\_\allowbreak{}dir\_\allowbreak{}up\_\allowbreak{}view}, \texttt{T\_\allowbreak{}diff\_\allowbreak{}up\_\allowbreak{}view}, \texttt{T\_\allowbreak{}total\_\allowbreak{}up\_\allowbreak{}view} & 아래 상향 투과 정의 \\
\texttt{TOA\_\allowbreak{}water\_\allowbreak{}signal\_\allowbreak{}I}, \texttt{T\_\allowbreak{}up\_\allowbreak{}rt\_\allowbreak{}valid} & 물 기원 TOA radiance와 상향 전달 계산 유효성 \\
\texttt{AOD\_\allowbreak{}ref} & 지정한 기준 파장에서의 에어로졸 광학두께 \\
\texttt{AOD\_\allowbreak{}ref\_\allowbreak{}nm} & AOD 기준 파장, 555 또는 865 nm \\
\texttt{AOD\_\allowbreak{}band} & \texttt{.\allowbreak{}mie} 소광비를 통해 계산한 실행 파장의 AOD \\
\texttt{AOD\_\allowbreak{}ext\_\allowbreak{}ratio} & 실행 파장 소광/기준 파장 소광; \texttt{AOD\_\allowbreak{}band=\allowbreak{}AOD\_\allowbreak{}ref*AOD\_\allowbreak{}ext\_\allowbreak{}ratio} \\
\texttt{sza}, \texttt{vza}, \texttt{raa}, \texttt{wl} & 실행 기하와 파장; RAA는 OCRT 규약 \\
\texttt{orders} & 비해양은 대기 Fourier 모드의 최대 사용 SOS 차수; 해양은 수중 풀이의 최대 사용 차수 \\
\texttt{conv} & 비해양은 대기 모드의 수렴 AND, 해양은 수중 풀이의 \texttt{all\_\allowbreak{}converged} 값 \\
\end{longtable}
}

\texttt{conv=\allowbreak{}1}은 코드의 해당 수렴 기준 통과를 뜻한다. 공간·각도 격자 독립성, 모델 간 정확도, 해양 전체 대기/수중 모든 경로의 독립 검증을 대신하지 않는다.

\subsection{비해양에 추가되는 키}
\label{sec:auto:05_output_formats:4.2}

{\TableFont
\begin{longtable}{@{}>{\raggedright\arraybackslash}p{\dimexpr 0.3400\linewidth-4.00pt\relax}>{\raggedright\arraybackslash}p{\dimexpr 0.6600\linewidth-4.00pt\relax}@{}}
\toprule
\textbf{키} & \textbf{의미} \\
\midrule
\endfirsthead
\toprule
\textbf{키} & \textbf{의미} \\
\midrule
\endhead
\bottomrule
\endfoot
\texttt{tau\_\allowbreak{}R} & 계산에 사용한 전체 Rayleigh 수직 광학두께 \\
\texttt{diag\_\allowbreak{}sunglint\_\allowbreak{}up\_\allowbreak{}flux\_\allowbreak{}ratio} & 직접 글린트 \texttt{rho\_\allowbreak{}I*cos(SZA)} 진단 \\
\texttt{diag\_\allowbreak{}TOA\_\allowbreak{}up\_\allowbreak{}flux\_\allowbreak{}ratio} & 기본 TOA \texttt{rho\_\allowbreak{}I*cos(SZA)} 진단 \\
\end{longtable}
}

비해양은 물에서 주어진 상향 출발광을 별도로 전파한 계산이 아니므로 \texttt{T\_\allowbreak{}diff\_\allowbreak{}up\_\allowbreak{}view}, \texttt{T\_\allowbreak{}total\_\allowbreak{}up\_\allowbreak{}view}, \texttt{TOA\_\allowbreak{}water\_\allowbreak{}signal\_\allowbreak{}I}가 \texttt{nan}, \texttt{T\_\allowbreak{}up\_\allowbreak{}rt\_\allowbreak{}valid=\allowbreak{}0}인 것이 정상이다.

\subsection{해양에 추가되는 키}
\label{sec:auto:05_output_formats:4.3}

{\TableFont
\begin{longtable}{@{}>{\raggedright\arraybackslash}p{\dimexpr 0.3400\linewidth-4.00pt\relax}>{\raggedright\arraybackslash}p{\dimexpr 0.6600\linewidth-4.00pt\relax}@{}}
\toprule
\textbf{키} & \textbf{의미} \\
\midrule
\endfirsthead
\toprule
\textbf{키} & \textbf{의미} \\
\midrule
\endhead
\bottomrule
\endfoot
\texttt{TOA\_\allowbreak{}rho\_\allowbreak{}atm\_\allowbreak{}I/\allowbreak{}Q/\allowbreak{}U}, \texttt{TOA\_\allowbreak{}rho\_\allowbreak{}water\_\allowbreak{}direct\_\allowbreak{}I/\allowbreak{}Q/\allowbreak{}U}, \texttt{TOA\_\allowbreak{}rho\_\allowbreak{}water\_\allowbreak{}sky\_\allowbreak{}I/\allowbreak{}Q/\allowbreak{}U}, \texttt{TOA\_\allowbreak{}rho\_\allowbreak{}water\_\allowbreak{}total\_\allowbreak{}I/\allowbreak{}Q/\allowbreak{}U}, \texttt{TOA\_\allowbreak{}rho\_\allowbreak{}glint\_\allowbreak{}on\_\allowbreak{}I/\allowbreak{}Q/\allowbreak{}U} & \ref{sec:auto:05_output_formats:3}절 TOA 성분 분해 \\
\texttt{Lu0plus}, \texttt{Qu0plus}, \texttt{Uu0plus} & 수면에서 공기로 나온 물 기원 Stokes 복사휘도 \\
\texttt{Ed0plus} & 공기 쪽 수면의 총 하향 조도 \\
\texttt{Ed0plus\_\allowbreak{}direct}, \texttt{Ed0plus\_\allowbreak{}diffuse} & 위 조도의 비산란 태양광/나머지 반구 적분 분해 \\
\texttt{Rrs0plus\_\allowbreak{}I}, \texttt{Rrs0plus\_\allowbreak{}Q}, \texttt{Rrs0plus\_\allowbreak{}U} & 위 복사휘도를 \texttt{Ed0plus}로 나눈 값 \\
\texttt{BRF0plus\_\allowbreak{}I}, \texttt{BRF0plus\_\allowbreak{}Q}, \texttt{BRF0plus\_\allowbreak{}U} & 각각 \texttt{\ensuremath{\pi}*Rrs0plus\_\allowbreak{}X} \\
\texttt{Lu0minus}, \texttt{Qu0minus}, \texttt{Uu0minus} & 수면 바로 아래 상향 Stokes 복사휘도 \\
\texttt{Ed0minus}, \texttt{Eu0minus} & 수면 바로 아래 총 하향/상향 조도 \\
\texttt{Ed0minus\_\allowbreak{}direct}, \texttt{Ed0minus\_\allowbreak{}diffuse} & 수중 하향 조도 분해 \\
\texttt{Eu0minus\_\allowbreak{}direct}, \texttt{Eu0minus\_\allowbreak{}diffuse} & 수중 상향 조도 분해; 현행 깊은 물/흑색 하단 모델에서 direct=0, diffuse=Eu0minus \\
\texttt{rrs0minus\_\allowbreak{}I}, \texttt{rrs0minus\_\allowbreak{}Q}, \texttt{rrs0minus\_\allowbreak{}U} & 수중 상향 Stokes를 \texttt{Ed0minus}로 나눈 값 \\
\texttt{BRF0minus\_\allowbreak{}I}, \texttt{BRF0minus\_\allowbreak{}Q}, \texttt{BRF0minus\_\allowbreak{}U} & 각각 \texttt{\ensuremath{\pi}*rrs0minus\_\allowbreak{}X} \\
\texttt{Kd0minus}, \texttt{Ku0minus} & 수면 부근 하향/상향 조도의 로그 감쇠 진단, m\textsuperscript{-}\textsuperscript{1} \\
\texttt{a\_\allowbreak{}w}, \texttt{b\_\allowbreak{}w}, \texttt{bb\_\allowbreak{}w}, \texttt{a\_\allowbreak{}dom}, \texttt{a\_\allowbreak{}chl}, \texttt{a\_\allowbreak{}pig}, \texttt{b\_\allowbreak{}pig}, \texttt{bb\_\allowbreak{}pig}, \texttt{a\_\allowbreak{}min}, \texttt{b\_\allowbreak{}min}, \texttt{bb\_\allowbreak{}min}, \texttt{a\_\allowbreak{}total}, \texttt{b\_\allowbreak{}total}, \texttt{bb\_\allowbreak{}total}, \texttt{omega\_\allowbreak{}total} & \ref{sec:auto:05_output_formats:7}절 IOP 대응표 \\
\end{longtable}
}

\texttt{Kd/\allowbreak{}Ku}는 수면과 첫 계산 깊이의 유한차분 \texttt{-ln(E(z1)/\allowbreak{}E(0))/\allowbreak{}z1}에서 얻는 수면 부근 진단이다. 특히 \texttt{Ku}는 표면 근처 조도 구배 때문에 작거나 음수일 수 있다. 이를 임의 깊이에서 일정한 고유 감쇠계수나 검증된 위성 Kd 제품으로 취급하지 않는다.

해양 단일 출력에는 \texttt{tau\_\allowbreak{}R} 키가 없고, 해양 격자에는 \texttt{Ku0minus}가 없다. 객체 내부 필드에 존재하는 값과 공개 파일에 실제 저장된 값을 구별한다.

\section{투과 관련 열: 무엇을 나눈 값인가}
\label{sec:auto:05_output_formats:5}

{\TableFont
\begin{longtable}{@{}>{\raggedright\arraybackslash}p{\dimexpr 0.3400\linewidth-4.00pt\relax}>{\raggedright\arraybackslash}p{\dimexpr 0.6600\linewidth-4.00pt\relax}@{}}
\toprule
\textbf{이름} & \textbf{정의·해석} \\
\midrule
\endfirsthead
\toprule
\textbf{이름} & \textbf{정의·해석} \\
\midrule
\endhead
\bottomrule
\endfoot
\texttt{T\_\allowbreak{}dir\_\allowbreak{}dn} / \texttt{direct\_\allowbreak{}transmittance} & 하향 비산란 직접광 감쇠. 평면평행에서는 \texttt{exp(-tau\_\allowbreak{}total/\allowbreak{}mu\_\allowbreak{}s)}; IPSS 사용 시 해당 직접광 구면 경로 적분을 적용 \\
\texttt{T\_\allowbreak{}diff\_\allowbreak{}dn\_\allowbreak{}hemi} / \texttt{diffuse\_\allowbreak{}dn\_\allowbreak{}hemi} & 수면 공기 쪽의 하향 확산 조도 / TOA 수평 입사량 \\
\texttt{T\_\allowbreak{}total\_\allowbreak{}dn\_\allowbreak{}hemi} / \texttt{irradiance\_\allowbreak{}transmittance} & \texttt{T\_\allowbreak{}dir\_\allowbreak{}dn+T\_\allowbreak{}diff\_\allowbreak{}dn\_\allowbreak{}hemi}; 수면 위 총 하향 조도 / TOA 수평 입사량 \\
\texttt{T\_\allowbreak{}diff\_\allowbreak{}dn\_\allowbreak{}dir} / \texttt{diffuse\_\allowbreak{}dn\_\allowbreak{}dir} & 요청 방향의 하향 확산 radiance로 정의한 방향 진단: \texttt{2*pi*mu\_\allowbreak{}v*L\_\allowbreak{}dn\_\allowbreak{}diff/\allowbreak{}(F0*mu\_\allowbreak{}s)} \\
\texttt{T\_\allowbreak{}dir\_\allowbreak{}up\_\allowbreak{}view} & 요청 VZA의 상향 비산란 감쇠 \texttt{exp(-tau\_\allowbreak{}total/\allowbreak{}mu\_\allowbreak{}v)} \\
\texttt{T\_\allowbreak{}total\_\allowbreak{}up\_\allowbreak{}view} & 실제 물 수출광의 각도 분포를 대기로 전달한 결과 \texttt{TOA\_\allowbreak{}water\_\allowbreak{}signal\_\allowbreak{}I/\allowbreak{}Lu(0+)} \\
\texttt{T\_\allowbreak{}diff\_\allowbreak{}up\_\allowbreak{}view} & \texttt{T\_\allowbreak{}total\_\allowbreak{}up\_\allowbreak{}view-T\_\allowbreak{}dir\_\allowbreak{}up\_\allowbreak{}view} \\
\texttt{TOA\_\allowbreak{}water\_\allowbreak{}signal\_\allowbreak{}I} & 물 신호만의 TOA 복사휘도, 상향 전달비의 분자 \\
\texttt{T\_\allowbreak{}up\_\allowbreak{}rt\_\allowbreak{}valid} & 필요한 물 기원 대기 bottom-source SOS 전파가 완료되었거나 대기가 없고 분모가 유효할 때 1; 그렇지 않으면 0 \\
\end{longtable}
}

\texttt{T\_\allowbreak{}diff\_\allowbreak{}dn\_\allowbreak{}dir}는 반구 적분값이 아니며 RAA 표본을 단순 합산한다고 \texttt{T\_\allowbreak{}diff\_\allowbreak{}dn\_\allowbreak{}hemi}가 되지 않는다. 방향과 적분 가중치를 포함해야 한다. \texttt{T\_\allowbreak{}total\_\allowbreak{}up\_\allowbreak{}view}는 그 장면의 물 수출광 분포에 의존한다. 물 IOP·기하에 독립적인 “대기만의 상향 투과 LUT”로 재사용하지 않는다. 1을 넘는 방향 전달비를 자동으로 잘라내지 않는다.

\texttt{T\_\allowbreak{}up\_\allowbreak{}rt\_\allowbreak{}valid=\allowbreak{}0}이면 상향 확산·총 전달비를 이용한 보정을 하지 않는다. 숫자가 필요한 경우 \texttt{T\_\allowbreak{}diff\_\allowbreak{}dn\_\allowbreak{}hemi}를 대신 넣는 상하 대칭 근사를 자동 적용해서도 안 된다.

다음 이름은 투과율이 아니다.


{\TableFont
\begin{longtable}{@{}>{\raggedright\arraybackslash}p{\dimexpr 0.3000\linewidth-5.33pt\relax}>{\raggedright\arraybackslash}p{\dimexpr 0.2300\linewidth-5.33pt\relax}>{\raggedright\arraybackslash}p{\dimexpr 0.4700\linewidth-5.33pt\relax}@{}}
\toprule
\textbf{LUT 열} & \textbf{단일 기하 키} & \textbf{정의} \\
\midrule
\endfirsthead
\toprule
\textbf{LUT 열} & \textbf{단일 기하 키} & \textbf{정의} \\
\midrule
\endhead
\bottomrule
\endfoot
\texttt{sunglint\_\allowbreak{}up\_\allowbreak{}flux\_\allowbreak{}ratio\_\allowbreak{}diag} & \texttt{diag\_\allowbreak{}sunglint\_\allowbreak{}up\_\allowbreak{}flux\_\allowbreak{}ratio} & \texttt{rho\_\allowbreak{}glint\_\allowbreak{}I*mu\_\allowbreak{}s} \\
\texttt{toa\_\allowbreak{}up\_\allowbreak{}flux\_\allowbreak{}ratio\_\allowbreak{}diag} & \texttt{diag\_\allowbreak{}TOA\_\allowbreak{}up\_\allowbreak{}flux\_\allowbreak{}ratio} & \texttt{rho\_\allowbreak{}TOA\_\allowbreak{}I*mu\_\allowbreak{}s} \\
\end{longtable}
}

이 값들은 이름에 flux가 있어도 반구 전체 상향 flux의 적분값이 아닌 \textbf{방향 radiance 기반의 과거 진단량}이다. \texttt{T\_\allowbreak{}total\_\allowbreak{}up\_\allowbreak{}view}로 해석하지 않는다.

\section{비해양 LUT CSV: 13개 열}
\label{sec:auto:05_output_formats:6}
실제 헤더는 다음과 같다.


\begin{ManualCode}
sza_deg,wavelength_nm,vza_deg,raa_deg,rho_I,rho_Q,rho_U,direct_transmittance,irradiance_transmittance,diffuse_dn_hemi,diffuse_dn_dir,sunglint_up_flux_ratio_diag,toa_up_flux_ratio_diag
\end{ManualCode}

앞의 네 열은 좌표, \texttt{rho\_\allowbreak{}I/\allowbreak{}Q/\allowbreak{}U}는 TOA 반사도다. 나머지는 \ref{sec:auto:05_output_formats:5}절 정의를 따른다. \texttt{direct\_\allowbreak{}transmittance}, \texttt{irradiance\_\allowbreak{}transmittance}, \texttt{diffuse\_\allowbreak{}dn\_\allowbreak{}hemi}는 같은 SZA·파장·대기·표면 케이스의 모든 방향에 반복되는 스칼라다. \texttt{rho}, \texttt{diffuse\_\allowbreak{}dn\_\allowbreak{}dir}, 두 \texttt{*\_\allowbreak{}diag}는 방향별 값이다.

이 CSV는 \texttt{0+} 쪽의 대기 계산을 나타낸다. \texttt{Rrs}, \texttt{rrs}, 물 IOP, 물 상향 전파 유효성을 제공하지 않는다. \texttt{rho}의 해석에는 파일에 없는 \texttt{surface}, \texttt{pressure}, AOD, \texttt{.\allowbreak{}mie}, 흡수 설정, 글린트 분리 설정도 필요하므로 실행 명령을 함께 보관한다.

\section{해양 full-grid CSV}
\label{sec:auto:05_output_formats:7}
실제 헤더는 아래 순서다.


\begin{ManualCode}
case_kind,sza_deg,wavelength_nm,wind_speed_ms,n_mu,mu_index,mu_view,vza_deg,raa_deg,Ed0plus_air,Ed0plus_direct,Ed0plus_diffuse,Lu0plus_I,Lu0plus_Q,Lu0plus_U,Rrs_I,Rrs_Q,Rrs_U,Ed0minus_water,Ed0minus_direct,Ed0minus_diffuse,Eu0minus_water,Eu0minus_direct,Eu0minus_diffuse,Lu0minus_I,Lu0minus_Q,Lu0minus_U,rrs_I,rrs_Q,rrs_U,TOA_rho_I,TOA_rho_Q,TOA_rho_U,T_dir_dn,T_diff_dn_hemi,T_total_dn_hemi,T_diff_dn_dir,T_dir_up_view,T_diff_up_view,T_total_up_view,TOA_water_signal_I,T_up_rt_valid,a_total,b_total,bb_total,omega_water,water_orders,water_converged,a_w,b_w,bb_w,a_chl,a_phyto_detritus,b_phyto_detritus,bb_phyto_detritus,a_dom,a_min,b_min,bb_min,Kd0minus
\end{ManualCode}

\subsection{좌표·식별·수렴}
\label{sec:auto:05_output_formats:7.1}

{\TableFont
\begin{longtable}{@{}>{\raggedright\arraybackslash}p{\dimexpr 0.3400\linewidth-4.00pt\relax}>{\raggedright\arraybackslash}p{\dimexpr 0.6600\linewidth-4.00pt\relax}@{}}
\toprule
\textbf{열} & \textbf{의미} \\
\midrule
\endfirsthead
\toprule
\textbf{열} & \textbf{의미} \\
\midrule
\endhead
\bottomrule
\endfoot
\texttt{case\_\allowbreak{}kind} & 고정 문자열 \texttt{ocean\_\allowbreak{}rrs\_\allowbreak{}grid} \\
\texttt{sza\_\allowbreak{}deg}, \texttt{wavelength\_\allowbreak{}nm}, \texttt{wind\_\allowbreak{}speed\_\allowbreak{}ms} & 케이스의 SZA, 파장, 풍속 \\
\texttt{n\_\allowbreak{}mu} & 대기 양의 Gauss 적분 노드 수; 출력 VZA 개수나 수중 노드 수가 아님 \\
\texttt{mu\_\allowbreak{}index} & 출력 VZA의 0부터 시작하는 행 묶음 인덱스; Gauss 노드 인덱스가 아님 \\
\texttt{mu\_\allowbreak{}view}, \texttt{vza\_\allowbreak{}deg}, \texttt{raa\_\allowbreak{}deg} & \texttt{cos(VZA)}, 공기 중 VZA, OCRT RAA \\
\texttt{water\_\allowbreak{}orders}, \texttt{water\_\allowbreak{}converged} & 수중 최대 SOS 차수와 수중 수렴 플래그 \\
\end{longtable}
}

\subsection{복사량·투과·반사도}
\label{sec:auto:05_output_formats:7.2}

{\TableFont
\begin{longtable}{@{}>{\raggedright\arraybackslash}p{\dimexpr 0.3400\linewidth-4.00pt\relax}>{\raggedright\arraybackslash}p{\dimexpr 0.6600\linewidth-4.00pt\relax}@{}}
\toprule
\textbf{열} & \textbf{단일 출력의 대응 키·정의} \\
\midrule
\endfirsthead
\toprule
\textbf{열} & \textbf{단일 출력의 대응 키·정의} \\
\midrule
\endhead
\bottomrule
\endfoot
\texttt{Ed0plus\_\allowbreak{}air} & \texttt{Ed0plus} \\
\texttt{Ed0plus\_\allowbreak{}direct}, \texttt{Ed0plus\_\allowbreak{}diffuse} & 같은 이름; 합은 \texttt{Ed0plus\_\allowbreak{}air} \\
\texttt{Lu0plus\_\allowbreak{}I}, \texttt{Lu0plus\_\allowbreak{}Q}, \texttt{Lu0plus\_\allowbreak{}U} & 각각 \texttt{Lu0plus}, \texttt{Qu0plus}, \texttt{Uu0plus} \\
\texttt{Rrs\_\allowbreak{}I}, \texttt{Rrs\_\allowbreak{}Q}, \texttt{Rrs\_\allowbreak{}U} & 각각 \texttt{Rrs0plus\_\allowbreak{}I/\allowbreak{}Q/\allowbreak{}U} \\
\texttt{Ed0minus\_\allowbreak{}water} & \texttt{Ed0minus} \\
\texttt{Ed0minus\_\allowbreak{}direct}, \texttt{Ed0minus\_\allowbreak{}diffuse} & 같은 이름; 합은 \texttt{Ed0minus\_\allowbreak{}water} \\
\texttt{Eu0minus\_\allowbreak{}water} & \texttt{Eu0minus} \\
\texttt{Eu0minus\_\allowbreak{}direct}, \texttt{Eu0minus\_\allowbreak{}diffuse} & 같은 이름; 합은 \texttt{Eu0minus\_\allowbreak{}water} \\
\texttt{Lu0minus\_\allowbreak{}I}, \texttt{Lu0minus\_\allowbreak{}Q}, \texttt{Lu0minus\_\allowbreak{}U} & 각각 \texttt{Lu0minus}, \texttt{Qu0minus}, \texttt{Uu0minus} \\
\texttt{rrs\_\allowbreak{}I}, \texttt{rrs\_\allowbreak{}Q}, \texttt{rrs\_\allowbreak{}U} & 각각 \texttt{rrs0minus\_\allowbreak{}I/\allowbreak{}Q/\allowbreak{}U} \\
\texttt{TOA\_\allowbreak{}rho\_\allowbreak{}I}, \texttt{TOA\_\allowbreak{}rho\_\allowbreak{}Q}, \texttt{TOA\_\allowbreak{}rho\_\allowbreak{}U} & 기본 해양 TOA 반사도 \\
\texttt{T\_\allowbreak{}dir\_\allowbreak{}dn}, \texttt{T\_\allowbreak{}diff\_\allowbreak{}dn\_\allowbreak{}hemi}, \texttt{T\_\allowbreak{}total\_\allowbreak{}dn\_\allowbreak{}hemi}, \texttt{T\_\allowbreak{}diff\_\allowbreak{}dn\_\allowbreak{}dir} & \ref{sec:auto:05_output_formats:5}절 하향 전달량 \\
\texttt{T\_\allowbreak{}dir\_\allowbreak{}up\_\allowbreak{}view}, \texttt{T\_\allowbreak{}diff\_\allowbreak{}up\_\allowbreak{}view}, \texttt{T\_\allowbreak{}total\_\allowbreak{}up\_\allowbreak{}view} & \ref{sec:auto:05_output_formats:5}절 상향 전달량 \\
\texttt{TOA\_\allowbreak{}water\_\allowbreak{}signal\_\allowbreak{}I}, \texttt{T\_\allowbreak{}up\_\allowbreak{}rt\_\allowbreak{}valid} & \ref{sec:auto:05_output_formats:5}절 상향 물 신호와 유효성 \\
\texttt{Kd0minus} & 수면 부근 하향 조도 감쇠 진단 \\
\end{longtable}
}

\subsection{IOP 메타데이터}
\label{sec:auto:05_output_formats:7.3}
계수는 계산에 \textbf{실제로 사용된} 값이다. 직접 IOP 입력 모드에서 성분별 값의 생물학적 분해를 추정하지 않는다. 총량에는 \texttt{a\_\allowbreak{}total}, \texttt{b\_\allowbreak{}total}, \texttt{bb\_\allowbreak{}total}을 사용한다.


{\TableFont
\begin{longtable}{@{}>{\raggedright\arraybackslash}p{\dimexpr 0.3000\linewidth-5.33pt\relax}>{\raggedright\arraybackslash}p{\dimexpr 0.2300\linewidth-5.33pt\relax}>{\raggedright\arraybackslash}p{\dimexpr 0.4700\linewidth-5.33pt\relax}@{}}
\toprule
\textbf{격자 열} & \textbf{단일 기하 키} & \textbf{의미} \\
\midrule
\endfirsthead
\toprule
\textbf{격자 열} & \textbf{단일 기하 키} & \textbf{의미} \\
\midrule
\endhead
\bottomrule
\endfoot
\texttt{a\_\allowbreak{}total}, \texttt{b\_\allowbreak{}total}, \texttt{bb\_\allowbreak{}total} & 동일 & 총 흡수·산란·후방산란, m\textsuperscript{-}\textsuperscript{1} \\
\texttt{omega\_\allowbreak{}water} & \texttt{omega\_\allowbreak{}total} & 수중 단일산란 알베도 \texttt{b\_\allowbreak{}total/\allowbreak{}(a\_\allowbreak{}total+b\_\allowbreak{}total)} \\
\texttt{a\_\allowbreak{}w}, \texttt{b\_\allowbreak{}w}, \texttt{bb\_\allowbreak{}w} & 동일 & 순수 해수 성분 \\
\texttt{a\_\allowbreak{}chl} & 동일 & 식물플랑크톤 흡수만의 성분; chlorophyll 농도 자체가 아님 \\
\texttt{a\_\allowbreak{}phyto\_\allowbreak{}detritus}, \texttt{b\_\allowbreak{}phyto\_\allowbreak{}detritus}, \texttt{bb\_\allowbreak{}phyto\_\allowbreak{}detritus} & \texttt{a\_\allowbreak{}pig}, \texttt{b\_\allowbreak{}pig}, \texttt{bb\_\allowbreak{}pig} & 코드에서 사용하는 식물플랑크톤·유기 쇄설 입자 성분 \\
\texttt{a\_\allowbreak{}dom} & 동일 & 용존 유기물 흡수 성분 \\
\texttt{a\_\allowbreak{}min}, \texttt{b\_\allowbreak{}min}, \texttt{bb\_\allowbreak{}min} & 동일 & 광물성 입자 성분 \\
\end{longtable}
}

\texttt{a\_\allowbreak{}chl}은 \texttt{a\_\allowbreak{}phyto\_\allowbreak{}detritus}의 식물플랑크톤 흡수 부분에 해당하므로 전체 흡수 합산 때 둘을 독립 성분처럼 다시 더하지 않는다. 모델마다 성분 경로가 다르므로 재구성 합계보다 출력 총량을 기준으로 검증한다.

\section{지정 기하 배치 CSV}
\label{sec:auto:05_output_formats:8}
현재 기본 헤더는 다음과 같다.


\begin{ManualCode}
case_id,sza_deg,vza_deg,raa_deg,wavelength_nm,rho_I_v2,rho_Q_v2,rho_U_v2
\end{ManualCode}

에어로졸 9열 입력 레이아웃이면 \texttt{wavelength\_\allowbreak{}nm} 뒤에 \texttt{aod\_\allowbreak{}ref}가 추가된다. 이 열은 기준 파장의 AOD이며 실행 파장의 AOD가 아니다. 기준 파장이 555인지 865인지는 입력 헤더·명령과 함께 보관한다. \texttt{\_\allowbreak{}v2}는 기존 파일 호환을 위한 열 이름이며 독립적인 물리 정의나 실행 파일 버전 번호가 아니다.

\texttt{-\allowbreak{}-\allowbreak{}output-\allowbreak{}advanced}를 주면 앞의 좌표 뒤에 다음 열을 쓴다.


\begin{ManualCode}
rho_I_v2,rho_I_ref,delta_abs,delta_rel_pct,delta_rel_abs_pct,tau_R_v2,tau_R_ref,n_orders_m0,n_orders_m1,n_orders_m2,walltime_ms
\end{ManualCode}

참조 Q/U가 없으면 마지막에 \texttt{rho\_\allowbreak{}Q\_\allowbreak{}v2,\allowbreak{}rho\_\allowbreak{}U\_\allowbreak{}v2}가 붙는다. 참조 Q/U가 모두 있으면 대신 다음 열이 붙는다.


\begin{ManualCode}
rho_Q_v2,rho_Q_ref,delta_abs_Q,delta_rel_pct_Q,delta_rel_abs_pct_Q,rho_U_v2,rho_U_ref,delta_abs_U,delta_rel_pct_U,delta_rel_abs_pct_U
\end{ManualCode}


{\TableFont
\begin{longtable}{@{}>{\raggedright\arraybackslash}p{\dimexpr 0.3400\linewidth-4.00pt\relax}>{\raggedright\arraybackslash}p{\dimexpr 0.6600\linewidth-4.00pt\relax}@{}}
\toprule
\textbf{열} & \textbf{정의·주의} \\
\midrule
\endfirsthead
\toprule
\textbf{열} & \textbf{정의·주의} \\
\midrule
\endhead
\bottomrule
\endfoot
\texttt{case\_\allowbreak{}id} & 입력 정수 식별자; 행별 격자 파일명에도 사용 \\
\texttt{rho\_\allowbreak{}I\_\allowbreak{}v2}, \texttt{rho\_\allowbreak{}Q\_\allowbreak{}v2}, \texttt{rho\_\allowbreak{}U\_\allowbreak{}v2} & 계산 TOA 반사도 \\
\texttt{rho\_\allowbreak{}I\_\allowbreak{}ref}, \texttt{rho\_\allowbreak{}Q\_\allowbreak{}ref}, \texttt{rho\_\allowbreak{}U\_\allowbreak{}ref} & 입력 참조 반사도 \\
\texttt{delta\_\allowbreak{}abs} & \texttt{rho\_\allowbreak{}I\_\allowbreak{}v2-rho\_\allowbreak{}I\_\allowbreak{}ref}; 이름에 abs가 있어도 \textbf{부호 있는 차이} \\
\texttt{delta\_\allowbreak{}rel\_\allowbreak{}pct} & \texttt{100*(rho\_\allowbreak{}I\_\allowbreak{}v2-rho\_\allowbreak{}I\_\allowbreak{}ref)/\allowbreak{}rho\_\allowbreak{}I\_\allowbreak{}ref} \\
\texttt{delta\_\allowbreak{}rel\_\allowbreak{}abs\_\allowbreak{}pct} & \texttt{abs(delta\_\allowbreak{}rel\_\allowbreak{}pct)} \\
\texttt{delta\_\allowbreak{}abs\_\allowbreak{}Q}, \texttt{delta\_\allowbreak{}abs\_\allowbreak{}U} & 각각 계산 Q/U - 참조 Q/U, 부호 있는 차이 \\
\texttt{delta\_\allowbreak{}rel\_\allowbreak{}pct\_\allowbreak{}Q}, \texttt{delta\_\allowbreak{}rel\_\allowbreak{}pct\_\allowbreak{}U} & \texttt{100*delta\_\allowbreak{}Q\_\allowbreak{}or\_\allowbreak{}U/\allowbreak{}abs(rho\_\allowbreak{}I\_\allowbreak{}ref)}; 분모가 Q/U 참조가 아니라 \textbf{참조 I의 절댓값} \\
\texttt{delta\_\allowbreak{}rel\_\allowbreak{}abs\_\allowbreak{}pct\_\allowbreak{}Q}, \texttt{delta\_\allowbreak{}rel\_\allowbreak{}abs\_\allowbreak{}pct\_\allowbreak{}U} & 위 백분율 차이의 절댓값 \\
\texttt{tau\_\allowbreak{}R\_\allowbreak{}v2}, \texttt{tau\_\allowbreak{}R\_\allowbreak{}ref} & 계산 및 입력 Rayleigh 광학두께 \\
\texttt{n\_\allowbreak{}orders\_\allowbreak{}m0}, \texttt{n\_\allowbreak{}orders\_\allowbreak{}m1}, \texttt{n\_\allowbreak{}orders\_\allowbreak{}m2} & Fourier 0/1/2 모드의 SOS 차수만 기록; 높은 모드의 최대 차수를 대표하지 않음 \\
\texttt{walltime\_\allowbreak{}ms} & 행 계산·관련 처리의 경과시간, 밀리초; 병렬 배치 전체 시간과 단순 합계가 같지 않음 \\
\end{longtable}
}

참조 I가 없으면 참조·오차 열은 비어 있다. 참조 I가 0이면 이 배치 출력기는 상대오차를 0으로 기록한다. 이를 “오차 없음”으로 판정하지 말고 상대오차 미정의로 다시 처리한다. 실패한 행의 반환코드는 별도 CSV 열에 없으므로 프로세스 종료코드, stderr, 값의 유한성도 확인한다.

이 \texttt{-\allowbreak{}-\allowbreak{}batch}는 \texttt{rt\_\allowbreak{}io\_\allowbreak{}run\_\allowbreak{}batch()}의 대기 지정 기하 경로다. 해양 구성성분·IOP의 일반 배치 인터페이스로 사용하지 않는다. 현행 \texttt{main.\allowbreak{}c}에서는 기체 흡수 초기화보다 앞에서 이 경로로 반환하므로 기체 옵션이 실제로 반영되지 않는다. 케이스 템플릿 복사도 제한되어 대기압의 연속값 등 모든 CLI 설정을 행에 전달하는 경로가 아니다. 기체·압력별 과학 LUT는 단일 케이스 격자를 반복 실행하고, 해양 작업은 \texttt{-\allowbreak{}-\allowbreak{}batch-\allowbreak{}full-\allowbreak{}grid} 또는 단일 기하 반복을 사용한다.

\texttt{case\_\allowbreak{}<id>.\allowbreak{}csv}를 함께 출력하는 경우 파일 내용은 \ref{sec:auto:05_output_formats:6}절의 비해양 13열 포맷이다. 같은 \texttt{case\_\allowbreak{}id}가 반복되면 같은 파일명을 쓰므로 입력 식별자를 유일하게 만든다. 행별 LUT 파일 열기 실패가 요약 성공 여부만으로 완전히 드러나지 않을 수 있으므로 파일 개수와 행 개수도 확인한다.

\section{CCRR 비교 CSV}
\label{sec:auto:05_output_formats:9}
이 포맷은 일반 운영용 시뮬레이션 파일보다 참조 비교를 위한 것이다. 헤더는 다음 묶음으로 구성된다. 묶음 안 이름은 실제 열 이름이다.


{\TableFont
\begin{longtable}{@{}>{\raggedright\arraybackslash}p{\dimexpr 0.3400\linewidth-4.00pt\relax}>{\raggedright\arraybackslash}p{\dimexpr 0.6600\linewidth-4.00pt\relax}@{}}
\toprule
\textbf{열 묶음} & \textbf{정의} \\
\midrule
\endfirsthead
\toprule
\textbf{열 묶음} & \textbf{정의} \\
\midrule
\endhead
\bottomrule
\endfoot
\texttt{case\_\allowbreak{}id,\allowbreak{}wavelength\_\allowbreak{}nm,\allowbreak{}chl\_\allowbreak{}mg\_\allowbreak{}m3,\allowbreak{}min\_\allowbreak{}g\_\allowbreak{}m3,\allowbreak{}adom\_\allowbreak{}440\_\allowbreak{}m-\allowbreak{}1,\allowbreak{}adom\_\allowbreak{}s\_\allowbreak{}nm-\allowbreak{}1} & 입력 사례 ID, 파장, Chl(mg m\textsuperscript{-}\textsuperscript{3}), 광물 농도(g m\textsuperscript{-}\textsuperscript{3}), 440 nm 유기물 흡수(m\textsuperscript{-}\textsuperscript{1}), 스펙트럼 기울기(nm\textsuperscript{-}\textsuperscript{1}) \\
\texttt{sza\_\allowbreak{}air\_\allowbreak{}deg,\allowbreak{}vza\_\allowbreak{}air\_\allowbreak{}deg,\allowbreak{}raa\_\allowbreak{}ccrr\_\allowbreak{}deg,\allowbreak{}raa\_\allowbreak{}ocrt\_\allowbreak{}deg,\allowbreak{}raa\_\allowbreak{}mapping} & 공기 중 기하와 전후 RAA; mapping 문자열은 \texttt{180-\allowbreak{}minus-\allowbreak{}ccrr} \\
\texttt{wind\_\allowbreak{}speed\_\allowbreak{}ms,\allowbreak{}sunglint\_\allowbreak{}decoupled} & 비교 경로가 고정하는 풍속 0, 글린트 분리 1 \\
\texttt{OCRT\_\allowbreak{}Ed0minus,\allowbreak{}OCRT\_\allowbreak{}Lu0minus,\allowbreak{}OCRT\_\allowbreak{}rrs,\allowbreak{}OCRT\_\allowbreak{}Ed0plus,\allowbreak{}OCRT\_\allowbreak{}Lu0plus\_\allowbreak{}no\_\allowbreak{}sky,\allowbreak{}OCRT\_\allowbreak{}Rrs} & 계산 조도·복사휘도·반사도 \\
\texttt{OCRT\_\allowbreak{}a\_\allowbreak{}total,\allowbreak{}OCRT\_\allowbreak{}b\_\allowbreak{}total,\allowbreak{}OCRT\_\allowbreak{}bb\_\allowbreak{}total} & 계산 총 IOP \\
\texttt{OCRT\_\allowbreak{}a\_\allowbreak{}w,\allowbreak{}OCRT\_\allowbreak{}b\_\allowbreak{}w,\allowbreak{}OCRT\_\allowbreak{}bb\_\allowbreak{}w,\allowbreak{}OCRT\_\allowbreak{}a\_\allowbreak{}dom,\allowbreak{}OCRT\_\allowbreak{}a\_\allowbreak{}pig,\allowbreak{}OCRT\_\allowbreak{}b\_\allowbreak{}pig,\allowbreak{}OCRT\_\allowbreak{}bb\_\allowbreak{}pig,\allowbreak{}OCRT\_\allowbreak{}a\_\allowbreak{}min,\allowbreak{}OCRT\_\allowbreak{}b\_\allowbreak{}min,\allowbreak{}OCRT\_\allowbreak{}bb\_\allowbreak{}min} & 계산 성분별 IOP; \ref{sec:auto:05_output_formats:7}절 의미와 동일 \\
\texttt{OCRT\_\allowbreak{}orders,\allowbreak{}OCRT\_\allowbreak{}conv} & 수중 차수·수렴 플래그 \\
\texttt{ref\_\allowbreak{}Ed0minus,\allowbreak{}ref\_\allowbreak{}Lu0minus,\allowbreak{}ref\_\allowbreak{}rrs,\allowbreak{}ref\_\allowbreak{}Ed0plus,\allowbreak{}ref\_\allowbreak{}Lu0plus\_\allowbreak{}no\_\allowbreak{}sky,\allowbreak{}ref\_\allowbreak{}Rrs,\allowbreak{}ref\_\allowbreak{}a\_\allowbreak{}total,\allowbreak{}ref\_\allowbreak{}b\_\allowbreak{}total,\allowbreak{}ref\_\allowbreak{}bb\_\allowbreak{}total} & 입력 파일의 해당 참조값 \\
\texttt{rel\_\allowbreak{}Ed0minus\_\allowbreak{}pct,\allowbreak{}rel\_\allowbreak{}Lu0minus\_\allowbreak{}pct,\allowbreak{}rel\_\allowbreak{}rrs\_\allowbreak{}pct,\allowbreak{}rel\_\allowbreak{}Ed0plus\_\allowbreak{}pct,\allowbreak{}rel\_\allowbreak{}Lu0plus\_\allowbreak{}no\_\allowbreak{}sky\_\allowbreak{}pct,\allowbreak{}rel\_\allowbreak{}Rrs\_\allowbreak{}pct,\allowbreak{}rel\_\allowbreak{}a\_\allowbreak{}total\_\allowbreak{}pct,\allowbreak{}rel\_\allowbreak{}b\_\allowbreak{}total\_\allowbreak{}pct,\allowbreak{}rel\_\allowbreak{}bb\_\allowbreak{}total\_\allowbreak{}pct} & 각 \texttt{100*(OCRT/\allowbreak{}ref-1)}; 참조 누락·0·비유한값이면 \texttt{nan} \\
\end{longtable}
}

\texttt{OCRT\_\allowbreak{}Lu0plus\_\allowbreak{}no\_\allowbreak{}sky}는 실제로 \texttt{res.\allowbreak{}Lu\_\allowbreak{}0plus\_\allowbreak{}view}를 기록한다. 열 이름만 보고 대기 하향 확산광이 물속에 들어가 만든 성분까지 제거되었다고 해석하지 않는다. 비교 시 참조 파일의 \texttt{no\_\allowbreak{}sky}가 수면 반사 하늘광 배제를 뜻하는지, 수중 입사 확산광까지 제거한 값인지 맞춰야 한다.

이 비교 경로는 \texttt{F\_\allowbreak{}sun=\allowbreak{}1}을 고정한다. 참조 \texttt{Lu/\allowbreak{}Ed}도 같은 입사량으로 정규화되었는지 확인한다. 이 경로도 일반 기체 흡수 초기화보다 먼저 실행되므로 기체 옵션을 포함하는 범용 검증 배치로 간주하지 않는다.

\section{결과를 자료로 저장할 때의 최소 메타데이터}
\label{sec:auto:05_output_formats:10}
기본 CSV만으로는 모든 실행 조건을 복원할 수 없다. LUT나 학습용 모의자료에는 다음을 함께 보관한다.

\begin{itemize}
\item Git 커밋과 실행 파일 식별자, 실행 명령, 사용한 환경변수.
\item RAA 정의와 SAA/VAA 변환식, 0--360\textdegree{} 보존 여부, Stokes 기준면·U 부호 처리.
\item 표면 모드, 글린트 분리 여부, 물 모델과 모든 구성성분/IOP 입력.
\item 파장, 대기압, AOD 기준 파장·값, 에어로졸 \texttt{.\allowbreak{}mie}, 기체 흡수 조건, 수온·염분·풍속.
\item 사용한 광학 LUT/위상함수 파일, 수치 설정, 종료코드와 오류 로그.
\item 원시 \texttt{Lu/\allowbreak{}Ed}를 저장하는 경우 내부 \texttt{F\_\allowbreak{}sun} 스케일과 실제 입사광 변환 여부.
\end{itemize}

표본별로 \texttt{Rrs=\allowbreak{}Lu0plus\_\allowbreak{}I/\allowbreak{}Ed0plus\_\allowbreak{}air}, \texttt{rrs=\allowbreak{}Lu0minus\_\allowbreak{}I/\allowbreak{}Ed0minus\_\allowbreak{}water}, \texttt{T\_\allowbreak{}total\_\allowbreak{}dn\_\allowbreak{}hemi=\allowbreak{}T\_\allowbreak{}dir\_\allowbreak{}dn+T\_\allowbreak{}diff\_\allowbreak{}dn\_\allowbreak{}hemi} 및 조도 direct+diffuse 합을 확인할 수 있다. 이런 내부 일관성은 RAA 규약이나 외부 정답과의 정확도를 대신 검증하지 않는다.

\subsection{구현 확인 위치}
\label{sec:auto:05_output_formats:10.1}
\begin{itemize}
\item \href{https://github.com/brtnt/OCRT/blob/54861c68e35ec5aaa3938fad53e93dc2dfa7eb1d/code/OCRT_C/src/rt_io.c}{단일 기하와 배치 출력기}: \texttt{rt\_\allowbreak{}io\_\allowbreak{}print\_\allowbreak{}single()}, \texttt{rt\_\allowbreak{}io\_\allowbreak{}run\_\allowbreak{}batch()}.
\item \href{https://github.com/brtnt/OCRT/blob/54861c68e35ec5aaa3938fad53e93dc2dfa7eb1d/code/OCRT_C/src/main.c}{각도 격자·CCRR CSV 출력기}: \texttt{run\_\allowbreak{}ocean\_\allowbreak{}rrs\_\allowbreak{}full\_\allowbreak{}grid\_\allowbreak{}csv()}, \texttt{run\_\allowbreak{}ccrr\_\allowbreak{}compare\_\allowbreak{}csv()} 및 마지막 비해양 LUT 출력 경로.
\item \href{https://github.com/brtnt/OCRT/blob/54861c68e35ec5aaa3938fad53e93dc2dfa7eb1d/code/OCRT_C/src/rt_types.h}{결과 필드 정의}: \texttt{rt\_\allowbreak{}result\_\allowbreak{}t}.
\item \href{https://github.com/brtnt/OCRT/blob/54861c68e35ec5aaa3938fad53e93dc2dfa7eb1d/code/OCRT_C/src/rt_solver.c}{실제 해양 성분·투과·정규화 계산}.
\item \href{https://github.com/brtnt/OCRT/blob/54861c68e35ec5aaa3938fad53e93dc2dfa7eb1d/code/OCRT_C/src/rt_water_rt.c}{수중 Kd/Ku 유한차분}.
\end{itemize}

